\scanpage{20}
En effet, on a un \deleted{\unsure{can}} 2-foncteur
\diagram{\begin{tikzcd}
\uuline{\operatorname{Brd}} \arrow[r,dash] & \uuline{\mathcal D} \\
O \arrow[r,mapsto] & C_O
\end{tikzcd}}
qui est une 2-équivalence. Malheureusement, on n'a pas de 2-foncteur naturel en sens inverse.

On supposera pour la suite que \deleted{$O^{\unsure{\natural|\sharp}}=A$} \inserted{les $O$ \deleted{sont des ordres maximaux d'une algèbre d'Azumaya sur $K$, $A$ Dedekind.} \inserted{alg. d'Azumaya}}, i.e. $A\xrightarrow{\sim}\operatorname{End}(1_C)$.
En termes de $C$, on va définir une \uline{gerbe} sur $A$ de \uline{groupe} $G_m$. Pour ceci, pour tt schéma \deleted{\unsure{S}} $T$ sur $A$, on va définir la catégorie des \deleted{\unsure{représentations}} \deleted{\uline{\unsure{primaires}}} \inserted{banalisations} de $C$ sur $T$, notée \deleted{$\uline{\operatorname{Rep}}_T(C)$} \inserted{$\mathcal G:\uline{\unsure{\operatorname{Ban}|\operatorname{Banal}}}(C)(T)$}, comme la catégorie des foncteurs $A$-\deleted{\unsure{projectifs}} \deleted{\unsure{fid}} linéaires
\diagram{\begin{tikzcd}
C \arrow[r,dash,"\varphi"] & \operatorname{Mod}(\deleted{\unsure{E}}\mathcal O_T)
\end{tikzcd}}
qui ont la propriété suivante: Pour tt $X\in C$ (ou \uline{un} $X\in\operatorname{Ob}C$, c'est pareil) \deleted{il existe} \inserted{posant $O=\operatorname{End}_C(X)$,} $\varphi(X)$ considéré comme \deleted{\unsure{\varphi}} $O\otimes_A\mathcal O_T$-Module est \deleted{\unsure{loc\textsuperscript{t}} libre de rang $1$} \inserted{une \uline{banalisation} de $O$}.

On constate alors: \deleted{\unsure{successivement}.} \deleted{\unsure{Le}} \unsure{foncteurs}

\deleted{a) \inserted{$T\longmapsto$} $\uline{\operatorname{Rep}}_T(C)$ est une gerbe pour la top. fppf.}
\margin{plus loin}
\deleted{b)} $\varphi\longmapsto\varphi(X)$ est une équivalence de la catégorie fibrée \deleted{\unsure{$\operatorname{End}_C$}} \inserted{$T\longmapsto\uline{\unsure{\operatorname{Ban}|\operatorname{Rep}}}(C)(T)$} et de la catégorie fibrée des banalisations de $O=\operatorname{End}(X)$; c'est donc une \uline{gerbe} de lien $G_m$. D'ailleurs l'opération de $G_m$ sur cette gerbe ne dépend pas du choix de $X$.
\mspage{19}
